\thispagestyle{empty}

\noindent{\sffamily\bfseries\large\color{Ink} How this book works}
\grule
\vspace{6pt}
\noindent{\sffamily\bfseries\large\color{Ink} The terms used in this book}
\grule
\vspace{6pt}


\tcap{About nouns}

\begin{flattbl}{colspec={X[0.85,l] X[2.15,l]}}
  Case & The form a noun, article or pronoun takes according to its job in the sentence. German has four: Nominativ, Akkusativ, Dativ, Genitiv. \\
  Gender & Whether a noun is masculine, feminine or neuter. It decides which article and adjective ending you use. \\
  Determiner & The word in front of a noun: an article (\textit{der}, \textit{ein}), a possessive (\textit{mein}), or a word like \textit{dieser}. \\
  der-word & A determiner that takes the same endings as \textit{der}: \textit{dieser}, \textit{jeder}, \textit{welcher}, \textit{mancher}, \textit{alle}. \\
  ein-word & A determiner that follows \textit{ein}, which has no ending in three of its sixteen slots: \textit{kein} and all the possessives. \\
  Weak noun & A small group of masculine nouns taking \textit{-n} or \textit{-en} in every case except the nominative singular, such as \textit{der Student}. \\
\end{flattbl}

\tcap{About adjectives}

\begin{flattbl}{colspec={X[0.85,l] X[2.15,l]}}
  Attributive & An adjective standing in front of a noun. Only these take endings. \\
  Predicative & An adjective standing after \textit{sein}, \textit{werden} or \textit{bleiben}. These never take an ending. \\
  Weak, mixed, strong & The three sets of adjective endings, chosen by what precedes the adjective: a der-word, an ein-word, or nothing. \\
\end{flattbl}

\clearpage
\thispagestyle{empty}

\noindent{\sffamily\bfseries\large\color{Ink} The terms used in this book \textnormal{\small\color{Faint}(continued)}}
\grule
\vspace{6pt}

\tcap{About verbs}

\begin{flattbl}{colspec={X[0.85,l] X[2.15,l]}}
  Weak verb & A verb whose past forms are made by adding \textit{-t-} to an unchanged stem: \textit{spielen, spielte, gespielt}. Sometimes called regular. \\
  Strong verb & A verb whose past forms are made by changing the stem vowel rather than adding \textit{-t-}: \textit{singen, sang, gesungen}. Around 200 exist; the 74 most common are listed in Part XV. \\
  Mixed verb & A small group that changes the vowel like a strong verb but takes the weak endings: \textit{bringen, brachte, gebracht}. \\
  Stem & The infinitive with \textit{-en} removed. Endings attach to it. \\
  Auxiliary & The helper verb that carries the person ending in a compound tense: \textit{haben}, \textit{sein} or \textit{werden}. \\
  Partizip II & The past participle, the form used with an auxiliary: \textit{gespielt}, \textit{gesungen}. \\
  Präteritum & The simple past, one word: \textit{ich spielte}. Mostly written German. \\
  Perfekt & The compound past, two words: \textit{ich habe gespielt}. Mostly spoken German. \\
  Konjunktiv II & The form used for hypotheticals and polite requests: \textit{ich hätte}, \textit{ich würde}. \\
  Separable verb & A verb whose stressed prefix detaches and moves to the end of the clause: \textit{aufstehen} $\rightarrow$ \textit{ich stehe auf}. \\
\end{flattbl}
